\section*{Broader Impact}

\paragraph{Deployment context.} Kazakhstan sits on the northern Tien Shan, one of the most seismically active mountain belts on Earth \citep{abdrakhmatov1996tienshanGPS}. Almaty, the country's largest city (population roughly 2 million), lies near a major active fault \citep{amey2021almaty} and was severely damaged by major earthquakes in 1887 and 1911. The Kazakhstan seismic network was not designed for comprehensive earthquake surveillance but for nuclear test monitoring, and its data centre operates with limited staff and no deep learning tools in its processing pipeline \citep{mikhailova2019kndc}. State-of-the-art deep learning models like PhaseNet could substantially extend detection capability in this setting, but as our results demonstrate, deploying them without understanding their failure modes carries its own risks.

\paragraph{Stakeholders and power asymmetries.} Well-resourced national agencies such as the United States Geological Survey and the Japan Meteorological Agency operate thousands of stations and can validate machine learning tools against dense monitoring networks before trusting them. The Kazakhstan National Data Centre, with a small team and a network built for a different purpose, does not have this capacity. The people most affected by monitoring quality are those living in seismically hazardous areas who depend on timely earthquake information but have no role in how monitoring systems are configured. Those most exposed to seismic risk are the least able to judge whether the tools meant to protect them are reliable.

\paragraph{The safety-critical gap.} Our results (\cref{sec:results}) show that PhaseNet's P-wave detection rate collapses from 83.9\% to 35.1\% under distribution shift while P-wave AUROC remains virtually unchanged, so PhaseNet's built-in confidence scores provide no warning that performance has degraded. This matters because automated picks propagate through earthquake catalogues into probabilistic hazard assessments \citep{cornell1968psha}, which inform building codes, emergency planning, and public warnings. Without a detection mechanism, unreliable picks from an unvalidated deployment region enter these systems with no indication that they should be treated differently from validated ones.

\paragraph{What the framework offers.} The regime-aware deferral framework does not eliminate the performance gap, but it provides a structured way to direct scarce analyst time toward periods of greatest environmental uncertainty, shifting the operational question from ``can we afford human review?'' to ``when should we prioritise it?'' The framework addresses precision-side failures: of the picks PhaseNet makes, the combined score identifies when they were made in poorly-understood regimes and routes them for review. Extending the approach to detect missed arrivals, which are PhaseNet's dominant failure mode under distribution shift, remains an open problem. The framework was evaluated on this study's dataset of 190 picks, and testing at the scale of continuous national monitoring is needed before operational deployment.

\paragraph{Transferability.} The distribution shift we study is not unique to Kazakhstan. PhaseNet is increasingly accessible for deployment to regions outside its training distribution through standardised interfaces like SeisBench \citep{woollam2022seisbench}, and the five Central Asian countries along the Tien Shan belt (Kazakhstan, the Kyrgyz Republic, Tajikistan, Turkmenistan, and Uzbekistan) share the same seismic source zone with varying levels of network coverage and catalogue completeness limitations \citep{poggi2024harmonizing}. The framework requires two inputs to transfer: a regional earthquake catalogue for HMM regime triggers, and a small number of analyst-verified picks to initialise the Beta-Bernoulli model. Both exist in these countries through the ISC Reviewed Bulletin, where available, making the architecture deployable without retraining PhaseNet or collecting new labelled data.

\paragraph{Broader implications.} To our knowledge, this is the first systematic evaluation of a state-of-the-art deep learning phase picker in the Kazakhstan seismic setting, and the first to demonstrate that PhaseNet's confidence scores provide no usable warning of performance degradation in this region. The finding that a widely-deployed model can fail silently under distribution shift, with confidence scores that look identical to in-distribution performance, has practical consequences for any agency considering adopting these tools. The regime-aware deferral framework offers one response: rather than requiring the model to diagnose its own failures, it uses structured environmental context that is already available through existing earthquake catalogues. The underlying principle extends beyond seismology: wherever deployment conditions are structured and observable, affect model reliability systematically, and are not reflected in the model's own uncertainty, external context can provide the deferral signal that internal confidence cannot. Examples include medical imaging across hospitals with different equipment profiles, autonomous vehicles across geographies with different road conditions, and environmental monitoring systems deployed across ecosystems.
